\section{Visualization and Interaction}
\label{sec:visualization}

The network view renders entries as nodes on a 2D HTML Canvas using d3-force. Nodes are colored by \texttt{sort} (deterministic hash-to-color mapping) and sized by configurable metrics.

\paragraph{Visual encoding.}

\begin{description}
  \item[Size] 11 metrics: in/out-degree, betweenness, closeness, eigenvector centrality, PageRank, clustering coefficient, HITS authority/hub, core number, degree centrality.
  \item[Color] Sort-based (default), metric gradient, cluster membership, or skeleton override.
  \item[Cluster] 5 methods: Louvain, spectral, label propagation, Girvan--Newman, $k$-core.
\end{description}

A ``tightness'' slider adjusts intra-cluster attractive force for visual compactness.

\paragraph{Interaction.}

\begin{itemize}
  \item Click to select a node and inspect it in the Detail View.
  \item Drag to reposition; scroll to zoom; pan to navigate.
  \item \texttt{Escape} to deselect; \texttt{Cmd+1/2/3} to switch panels.
\end{itemize}

Selection propagates across all panels: the Detail View shows the entry's content, and MDX components update their color to reflect the current visualization settings.

\paragraph{Entry macros.}
The Read View renders \texttt{.mdx} documentation files that can embed astrolabe entries inline via two \LaTeX-style macros:

\begin{itemize}
  \item \verb|\entryref{hash}{display text}| --- an inline clickable link to an entry. Rendered as colored text; clicking selects the entry in the network view.
  \item \verb|\entryblock{hash}| --- a block-level display of an entry. Fetches the entry's record from the backend and renders its sort label, title, and notes with a colored left border.
  \item \verb|\entryblock{hash}{collapsible}| --- same as above but collapsed by default; click to expand.
  \item Nesting: \verb|\entryblock{hash}{\entryblock{child}{collapsible}}| --- an entry block containing a collapsible child block, enabling hierarchical presentation.
\end{itemize}

These macros are preprocessed before Markdown rendering: a parser matches \verb|\entryref{...}{...}| and \verb|\entryblock{...}| patterns (with recursive brace matching for nested blocks) and converts them to HTML elements with \texttt{data-entry} attributes. React components then fetch the entry data and render with sort-derived colors. This allows documentation to be a living view of the knowledge network: entry content updates automatically when the underlying data changes, and colors stay synchronized with the network visualization settings.
